\begin{textsample}{POS Dim 3 – mg – Score 16.00 – mg/mg\_inf\_avaliacao.txt}  \label{ex:f3_pos_265}
2.9 Avaliação e Registro na Educação \textbf{Infantil} A avaliação da aprendizagem anda de \textbf{mãos} dadas com a ação de planejar.
Segundo Blasis et al. ( 2013 ), a avaliação é um processo crítico e permanente da prática pedagógica, que constitui o ponto de partida para o aperfeiçoamento desta prática.
Neste sentido, a avaliação da aprendizagem é concebida para o diagnóstico das necessidades pedagógicas e possibilita a identificação das condições formativas ideais para a \textbf{criança} em seu contexto.
Ou seja, ela informa sobre os processos de aprendizagem da \textbf{criança} e subsidia o trabalho do professor, apontando as necessidades de continuidade, de avanços ou de mudanças no seu planejamento e no desenvolvimento das ações educativas.
Segundo Luckesi ( 1996, p. 81 ), “ a avaliação deverá ser assumida como um instrumento de compreensão do estágio de aprendizagem em que se encontra o estudante, tendo em vista tomar decisões suficientes e satisfatórias para que possa avançar no seu processo de aprendizagem ”.
Ou seja, quanto mais o professor conhecer as formas pelas quais as \textbf{crianças} aprendem, melhor será sua intervenção pedagógica.
Assim, a avaliação torna -se o instrumento de mediação entre as propostas organizadas pelo professor e as aprendizagens da \textbf{criança}, é o fio de comunicação entre as formas de ensinar e formas de aprender.
Esta concepção de avaliação evidencia que a função da avaliação na Educação \textbf{Infantil}, como também em todos os níveis escolares, não é classificatória, não é medida e não deve ser realizada ao final do processo.
Ou seja, a avaliação está a serviço da regulação do ensino e da aprendizagem, portanto, ela deve ser sistemática e contínua, tendo como objetivo principal a melhoria da ação educativa, conforme ressalta Jussara Hoffmann ( 2012, p. 25 ) : > Um dos pressupostos básicos dessa prática é, justamente, o seu caráter investigativo e mediador, não constatativo.
A permanente \textbf{curiosidade} dos professores sobre as \textbf{crianças} é premissa básica da avaliação em Educação \textbf{Infantil}, e não a intenção de julgar como positivo ou negativo o que uma \textbf{criança} é ou não capaz de fazer e de aprender.
Essa afirmação da autora nos remete à compreensão de que a avaliação na Educação \textbf{Infantil} tem como foco fornecer informações acerca dos processos e das aprendizagens das \textbf{crianças}, atendendo aos princípios de que as \textbf{crianças} aprendem de formas diferentes, em tempos diferentes, a partir de diferentes vivências pessoais e experiências anteriores.
Portanto, o papel da escola é desenvolver o potencial e as habilidades de cada indivíduo sem qualquer tipo de discriminação, oportunizando a sua participação de forma ativa e abrangente nos contextos sociais e culturais.
De acordo com a Lei de Diretrizes e Bases da Educação ( 1996 ), o artigo 31 estabelece que “ [... ] a avaliação far -se-á mediante o \textbf{acompanhamento} e o \textbf{registro} do seu desenvolvimento, sem o objetivo de promoção, mesmo para o acesso ao ensino fundamental ”.
Nestas condições, fica estabelecido que é de responsabilidade das \textbf{instituições} de ensino de Educação \textbf{Infantil} a criação de procedimentos para \textbf{acompanhamento} do trabalho pedagógico e para a avaliação do desenvolvimento das \textbf{crianças}, utilizando a observação crítica, criativa e contínua, além de instrumentos variados de \textbf{registros} e de documentação.
A observação e o \textbf{registro} constituem -se nos principais instrumentos de que o professor da Educação \textbf{Infantil} dispõe para \textbf{apoiar} sua prática avaliativa.
Por meio da observação sistemática, o professor pode registrar contextualmente os processos de aprendizagens das \textbf{crianças}, a qualidade das interações e acompanhar os processos de desenvolvimento, obtendo informações sobre as experiências das \textbf{crianças} na \textbf{instituição}.
Esta observação e as diferentes formas de \textbf{registro} ( \textbf{diário} de campo, \textbf{livro} de memória do grupo, portfólios, fotografias, vídeos, depoimento de pais etc. ) fornecem aos professores uma visão integral das \textbf{crianças}, ao mesmo tempo em que revelam suas particularidades. ( BRASIL, 1998 ). > Para acompanhar e avaliar as \textbf{crianças}, é importante a observação sistemática, crítica e criativa do comportamento de cada uma, dos grupos, das \textbf{brincadeiras} e interações entre elas no cotidiano.
Além disso, é fundamental a utilização de múltiplos \textbf{registros} realizados por \textbf{adultos} e \textbf{crianças} ( relatórios, fotografias, desenhos, álbuns etc. ), em diversificados momentos.
Conhecer as preferências das \textbf{crianças}, a forma de elas participarem nas atividades, seus \textbf{parceiros} prediletos para a realização de diferentes tipos de tarefas, suas narrativas podem ajudar o/a professor/a a reorganizar as atividades de modo mais adequado, ao alcance dos propósitos \textbf{infantis} e das aprendizagens coletivamente trabalhadas [... ] ( BRASIL, 2009, p. 17 ).
Neste sentido, entende -se que quanto mais variada for a escolha de instrumentos de \textbf{registro} do professor, mais fontes de informação ele terá para apoiá -lo na sua reflexão, construção de sentidos e criação de novos caminhos no processo de ensino-aprendizagem.
Conclui -se que a avaliação, em aspectos gerais, baseia -se nas experiências vividas pela \textbf{criança}, na observação e utilização de múltiplos \textbf{registros} dessas experiências, que visam administrar, de forma contínua, os processos de desenvolvimento e de aprendizagem da \textbf{criança} na Educação \textbf{Infantil}.
Cabe mencionar ainda que, tendo em vista a busca por qualificação pedagógica, política e estrutural do serviço educativo que é ofertado, faz -se necessário propor caminhos para a avaliação da Educação \textbf{Infantil} que esteja implicada e comprometida com o desenvolvimento e aprendizagem de \textbf{meninos} e \textbf{meninas}.
Consequentemente, neste contexto, estão incluídos : a estrutura física dos locais de funcionamento, \textbf{brinquedos}, materiais de uso pedagógico disponíveis, a composição do acervo literário, o PPP da \textbf{instituição}, a formação docente e as possibilidades de aperfeiçoamento profissional em serviço, como também o relacionamento com as famílias e a comunidade em que está inserida.
De acordo com as DCNEI ( 2010 ), a \textbf{criança} constrói sua identidade pessoal e coletiva, na \textbf{percepção} de si mesma e do outro, nas \textbf{brincadeiras}, no faz de conta, nas experimentações, vivências, narrativas e investigações, entre tantas outras vivências.
Dessa maneira, a forma como essas experiências são planejadas, organizadas e avaliadas, considerando os espaços, tempos de permanência, materialidades e relações, podem se constituir em potenciais contextos de aprendizagem, expressos na forma como a \textbf{criança} aprende, interage e vivencia descobertas.

% matched lemmas: acompanhamento, adulto, apoiar, brincadeira, brinquedo, criança, curiosidade, diário, infantil, instituição, livro, menino, mão, parceiro, percepção, registro
\end{textsample}
